\documentclass[a4paper,11pt,titlepage]{article}
\parskip 3pt

%% %%%%%%%%%%%%%%%%%%%% BEGIN PACKAGES %%%%%%%%%%%%%%%%%%%%

%%\usepackage[top=1in, bottom=1in, left=1in, right=1in]{geometry}

%%\usepackage{fullpage}
\usepackage{url}
\usepackage{hyperref}
\usepackage{listings}
\hypersetup{pdfborder=0 0 0}

%% Images
\usepackage{graphicx}

%% Side-by-side images
\usepackage{subfig}

%% Wrapped figures
\usepackage{wrapfig}

%% Drawing trees and other stuff
\usepackage{tikz}
\usetikzlibrary{trees,arrows}
\usepackage{amssymb}

%% Pseudocode
\usepackage[noend]{algorithmic}
\algsetup{indent=1.5em}

%% Linux Libertine
\usepackage[T1]{fontenc}
\usepackage{libertine}
\renewcommand*\oldstylenums[1]{{\fontfamily{fxlj}\selectfont #1}}

\let\stdhref\href
\renewcommand{\href}[2]{\stdhref{#1}{\texttt{#2}}}

\newcommand{\mailto}[1]{\href{mailto:#1}{#1}}

\let\stdsection\section         % because LaTeX cannot handle
                                % recursive commands
\renewcommand{\section}{\newpage\stdsection}

\let\tikzsquare\square
\renewcommand{\square}{\ensuremath\tikzsquare}

\newcommand{\code}[1]{\texttt{#1}}
\newcommand{\buzz}[1]{\emph{#1}}
\newcommand{\magic}[1]{\buzz{#1}}

\newcommand{\myhref}[1]{\href{http://#1}{#1}}

%% Clubs and widows. Avoid them like the plague.
\clubpenalty = 10000
\widowpenalty = 10000
\displaywidowpenalty = 10000

%% %%%%%%%%%%%%%%%%%%%% END COMMANDS %%%%%%%%%%%%%%%%%%%%

%Visigoth is a tool to generate, analyse and visualise Small World
%Networks. These are particular kinds of graphs which \emph{look} like naturally
%grown networks and share many mathematical properties with them.  For instance,
%friendship graphs and graphs of followers in modern web-based ``social
%networks'' are Small World Networks.

%As an educational tool Visigoth allows users to explore Small World Networks and
%hence can be used as an introduction to them. By exploring these networks users
%will have a firmer, more intuitive grasp of the mathematical properties they
%exhibit.

%The networks may be generated using several algorithms, allowing users to
%compare their output and its adherence to the properties of Small World
%Networks. The network generation algorithms in Visigoth form a smooth
%progression in the number of exhibited desired properties from completely random
%graphs, generated by the Erd\H{o}s-R\'{e}nyi algorithm, to fully compliant Small
%World Networks, generated by the Preferential Attachment
%algorithm. Additionally, Visigoth can connect to social networking platforms
%such as \Twitter to fetch and display real-world networks.

%We should keep it 2 pages long at most and should have the same structure as 
%the Report


\begin{document}
\title{\includegraphics[width=2cm]{v.png}\\
  \Huge Visigoth\\\Large Graph visualisations}
\author{
  Andreea-Ingrid Funie\\\mailto{aif109@doc.ic.ac.uk}\and
  Alexandru Scvor\c tov\\\mailto{as10109@doc.ic.ac.uk}\and
  Francesco Mazzoli\\\mailto{fm2209@doc.ic.ac.uk}\and
  Marc-David Haubenstock\\\mailto{mh808@doc.ic.ac.uk}\and
  Maximilian Staudt\\\mailto{ms9109@doc.ic.ac.uk}
}
\date{January 2012}
\maketitle

\begin{description}
	\item Introduction
	% High Level, Nontechnical Description Why you should buy this product/listen to
	% this presentation? What is the functionality of the product?

	\item Technical Description
	%B. Short Technical Description; 
	%Short introduction into technologies used
	%Design of your software, possibly including a diagram of the major components
	% of the project; Main achievements

	\item Software Engineering Issues
	%C. Software Engineering Issues:
	%What technology was used and why; what other technology was considered but not
	% used and why; Any technical challenges encountered and how addressed;
	%Any risks anticipated, and how mitigated? Any collaboration/coordination 
	%difficulties encountered and how addressed; Development and testing methods
	% and/or tools used; comparison of plans with actual achievements; Estimates of 
	%length of code in each of the components, or any other comparable measure of
	% the effort required; Summary of each team member's contributions

	\item Validation
	%D. Validation and Conclusions How did you validate your product? 

	\item Conclusions
	%Was the project successful? What did you learn? What might you have done
	% differently?

\end{description}

\end{document}
